\section{Data protection and cybersecurity}
Seven-day retention follows the assumed incident-reporting cycle and capacity budget, not a universal legal obligation. The data controller justifies necessity; incident evidence has a separate purpose and deletion deadline. Employment use must also consider Article 22$^2$ of the Polish Labour Code, including monitoring purposes, information duties, and privacy in protected rooms. The model has no cameras in toilets, staff welfare rooms, or at SOC/office desks.\cite{L-GDPR,L-KP}

Before actual commissioning, the controller establishes a GDPR Article 6 basis and assesses the need for a DPIA under Article 35, considering scale, regularity, and public space. Zone signs, full Article 13 information, privacy masks, no audio, and export restrictions are assumed. C08 person detection creates no biometric identifier; facial or emotion recognition is not designed. Cards instead of biometrics reduce special-category data processing.\cite{L-GDPR}

Every person has an individual account. Roles separate viewing, export, card issuance, and administration. Unused ports are disabled; cameras and controllers have no Internet access. Administrative access is temporary and logged; the local computer has disk encryption and session locking. Firewall rules allow only described services between specified devices; disabled UPnP, P2P, and cloud features prevent undeclared channels. HTTPS, SMB3, and OSDP Secure Channel are used according to device support. A VLAN alone is not assumed to provide confidentiality.\cite{S-VMS,S-KD,S-FW}

Configuration is backed up after an approved change and stored on the administrator's encrypted medium outside the recording array. Keys and passwords are excluded from the model and repository. Updates require card, reader, and VMS compatibility testing and a restoration path. Fire-system software is not updated from an IT operator account. The project's 90-day log retention supports incident analysis and requires justification like other retention periods.
